% !TEX root = ../main.tex

\chapter{Related Work}
\label{ch:related-work}
Integrating task and motion planning with partial observability is challenging, particularly when a robot must act on incomplete information about its environment. This chapter reviews previous work along the two axes the thesis combines: spatial representations a planner can reason over (\cref{sec:rw-spatial}), the axis on which this thesis contributes, and planning approaches for TAMP under partial observability (\cref{sec:rw-potamp}).

\section{Spatial Representations for Planning under Occlusion}
\label{sec:rw-spatial}

The contribution of this thesis is a spatial representation, so the closest related work is not only other planners but other ways of carving up space. The standard choices, uniform voxel grids and octree occupancy maps such as OctoMap \cite{hornung2013octomap}, answer where matter is: each cell is occupied, free, or unknown, and unobserved space is represented implicitly, by the absence of evidence (\cref{sec:spatial_representation}). Nothing in the structure names which unknown volume matters to the task, which object hides it, or what must move for it to become visible; a planner searching for a hidden object would have to recover that meaning from raw geometry. Our evaluation keeps this family as the controlled baseline: the \texttt{uniform} variant (\cref{sec:baselines}) is exactly such a grid at the same cell size as the adaptive partition, so the measured gap isolates the structure of the partition rather than its resolution.

A second line abstracts space by task relevance instead of occupancy: learned critical regions \cite{molina2020learn, shah2022abstractions} partition the configuration space around regions that matter to solutions, such as narrow passages, and plan hierarchically over the resulting abstract states (\cref{sec:spatial_representation}). The regions, however, are predicted once per environment and held fixed, and the abstraction carries no notion of occlusion, so it neither names the volume hidden behind an object nor adapts when that object is moved; later work in this line that does adapt targets reinforcement learning rather than TAMP \cite{nayyar2025option}. Boxels keep the task-relevance idea but construct the regions geometrically, from the detected objects and their sightlines to the camera, and rebuild them whenever the scene changes.

Planning systems that do reason about occlusion represent it as a property of a probabilistic belief rather than as a region of space. TAMPURA's hidden-object benchmark keeps a visibility voxel grid that weights the success probability of its \texttt{look} action \cite{curtis2024partially}; IBSP carries maximum-likelihood occlusion fluents \cite{hadfield2015modular}; SS-Replan derives an occlusion predicate by ray-casting over its pose belief \cite{garrett2020online}; and occlusion-aware retrieval searches the likely poses of a hidden target under a learned distribution \cite{bejjani2021occlusion}. In each case occlusion is a check or a weight computed against a distribution over object poses, not a named occluded volume that a planner can sense, clear, or avoid re-blocking as symbolic state. That named-volume representation, constructed geometrically and held as deterministic planning state, is what the Boxel partition adds; \cref{sec:rw-potamp} reviews the planning frameworks it is evaluated against.

\section{TAMP under Partial Observability}
\label{sec:rw-potamp}
Existing approaches to TAMP under partial observability generally either use formal, probabilistic models or more ad-hoc, reactive methods, or a combination of both.

\subsection{POMDP-based TAMP Solutions}
Several works formulate planning under uncertainty as a POMDP to handle it in a principled, risk-aware way \cite{curtis2024partially, saleem2024pomdp}. \Ac{tampura} \cite{curtis2024partially}, for example, learns a coarse abstract model of each given controller's preconditions and effects, builds a non-deterministic MDP, and solves it with \ac{lao}, an uncertainty-aware solver \cite{hansen2001lao}, to produce risk-aware, information-gathering plans. \cite{saleem2024pomdp}, by contrast, addresses contact-rich \emph{manipulation} under object-pose uncertainty (not full task-and-motion planning), computing partial policies with an anytime heuristic-search POMDP solver.
However, a persistent challenge is the representation of belief over continuous properties, particularly object poses. TAMPURA \cite{curtis2024partially} samples object placements in continuous pose space and, in its real-robot pipeline, draws its object-pose belief from a perception front-end, Bayes3D \cite{gothoskar2023bayes3d}: a generative inverse-graphics model that infers a posterior over 3D scenes by rendering candidate object arrangements and scoring them against the observed depth image.

\subsection{Partially Grounded Plans in TAMP}
In contrast to formal probabilistic models, other TAMP systems handle uncertainty by deferring it to execution \cite{pan2024task}. These methods generate partially grounded plans, leaving ``gaps'' for actions with unknown outcomes. During execution, specialized closed-loop behaviors (hand-designed or learned) attempt to fill these gaps. If a controller fails, the failure is treated as a new constraint, and the system replans. This reactive approach avoids the need to model outcome probabilities explicitly. However, because it does not reason about uncertainty proactively, it can struggle with problems that require multi-step, strategic information gathering, such as deciding which of several actions is most likely to reveal a hidden object.

\subsection{Modern Approaches to TAMP under Uncertainty}
Recent works have leveraged learning and large models to address partial observability, offering different strategies than our own.

\cite{ma2025task} adds a strategic reasoning layer on top of a standard planner: it first reasons backward from an underspecified goal to decide which objects matter (Task-level Backward Search), then uses a constraint graph to order object moves collision-free through clutter (Object Manipulation Constraint Graph). We differ at the representation level rather than by adding a layer: we introduce a geometric belief representation, Boxels, so that the planner treats which regions are observed versus occluded as first-class planning state and information-gathering becomes a core planner action rather than a separate reasoning stage; the visibility signal is currently supplied by an oracle.

Another approach sources the planner's belief from large foundation models. \cite{zhao2025seeing} uses a \ac{vlm} as an uncertainty estimator: the VLM answers the planner's perceptual queries about an image (for example, ``Is the drawer empty?''), and the resulting belief feeds a \emph{symbolic} belief-space planner, a pairing of learned perception with structured reasoning that is reported to outperform end-to-end VLM planners. CoCo-TAMP~\cite{kim2026llmguided} is similar in spirit, drawing on the commonsense priors of a \ac{llm} to shape the belief over where task-relevant objects are likely to be, again within a partially observable TAMP loop. The difference from our work is the \emph{source} of the belief: theirs is learned (a VLM or LLM), whereas ours is a structured, object-centric geometric partition into discrete Boxels, with object poses currently supplied by a ground-truth oracle rather than a learned detector. Integrating a real detector is left to future work.

Finally, some methods move away from symbolic planning altogether and instead use end-to-end learning. \cite{bai2025learning} (TAVP), for example, uses a separate reinforcement-learning policy to select task-relevant \emph{virtual} camera viewpoints and re-render the scene from them, rather than physically moving the camera, and then produces the manipulation action with an end-to-end learned policy rather than a planner. This is fundamentally different from our model-based approach, where looking around is not a learned reflex but an explicit action chosen by the planner because its world model indicates that knowledge is missing. Because our method separates planning from execution, it is more transparent and can be retargeted to new goals without retraining a policy; the cost, however, is that a genuinely new goal requires authoring new PDDL actions and their streams, not just a new goal specification.

\subsection{Belief-Space Planning and Replanning}
A long line of work plans directly in \emph{belief space}, treating the robot's distribution over world states as the object of planning. \cite{kaelbling2013integrated} plans over the robot's belief using hierarchical goal regression, while \cite{hadfield2015modular} introduces a modular interface that exchanges information between an off-the-shelf classical planner and the geometric layer as the plan is built, determinising the problem under the maximum-likelihood-observation assumption so that a successful plan generates the observations it needs. A practical alternative to maintaining an explicit stochastic model is to \emph{determinise and replan}: SS-Replan \cite{garrett2020online} plans deterministically in a hybrid belief space and replans whenever a new observation arrives, including sensing actions within a single plan. Our system shares this determinise-and-replan strategy: we commit to an optimistic deterministic plan and replan on each observation rather than synthesising a probabilistic policy. What we add is separate from the replanning loop: a spatial belief in which occlusion is an explicit, named planning state, so information-gathering is explicitly composed within the planning representation.

\subsection{Partially Observable Deterministic Planning}
A separate tradition handles partial observability \emph{without} probabilities, by reasoning about what is \emph{known}. A prominent line within partially observable deterministic (POD) planning represents knowledge with explicit \emph{knowledge literals} (\cref{sec:pod}) and reduces the problem to classical planning by compilation; other POD planners instead search belief space directly, as in Contingent-FF~\cite{hoffmann2005contingent} and MBP~\cite{bertoli2001planning}. The translation-based contingent planner CLG \cite{albore2009translation} and the classical-replanning K-replanner \cite{bonet2011planning} established this reduction; LW1 \cite{bonet2014flexible} made it scalable with a \emph{linear} translation; and a parallel route compiles partial observability to \ac{fond} planning \cite{muise2014computing}. This thesis builds directly on the POD line: we carry Know-If Fluents (e.g.\ whether a region's contents are known) inside the PDDL state and let a classical planner reason over them.

To our knowledge, the deterministic, knowledge-literal POD line has not previously been combined with task and motion planning, nor used a geometric belief representation that makes \emph{which regions are occluded} part of the symbolic state. The systems that do reason about occlusion do so over probabilistic beliefs or learned models (\cref{sec:rw-spatial}), and partially observable TAMP systems increasingly source their belief from learned models \cite{zhao2025seeing, kim2026llmguided} or probabilistic abstractions \cite{curtis2024partially}; we are not aware of prior work that represents occluded volumes as named, object-centric symbolic state resolved by sensing within a PDDLStream TAMP loop, the gap this thesis addresses.

\section{Summary and Research Gap}
These approaches occupy distinct points in the design space. POMDP-based methods such as TAMPURA \cite{curtis2024partially} reason about uncertainty proactively but rely on a learned probabilistic transition model, and they keep the spatial belief sub-symbolic, so occluded regions are never a named, inspectable state that the planner can target. Reactive, partially-grounded methods \cite{pan2024task} avoid probabilistic modelling but cannot plan multi-step information-gathering. Belief-space replanning \cite{garrett2020online} determinises and replans, as we do, but does not make occlusion a first-class planning entity. The POD tradition \cite{bonet2014flexible} reasons about knowledge symbolically and deterministically, but has not been applied to TAMP or to spatial occlusion. We build on these lines: a POD task-and-motion planner whose spatial belief is an adaptive, object-centric partition (Boxels) in which occluded regions and their observed/occluded status are an explicit symbolic state. The planner thus reasons about \emph{where} information is missing and composes deliberate look-and-clear actions, without the cost and opacity of a learned probabilistic model. On the spatial side, \cref{sec:rw-spatial} positions the partition against the existing representations: occupancy structures carry no semantic notion of an occluded volume, learned task-relevance abstractions are predicted once per scene rather than reconstructed as objects move, and occlusion-aware planners keep occlusion inside a probabilistic belief. None names the occluded volumes that sensing must resolve, which is what the Boxel partition adds.

This position yields three advantages, each with a limitation we state plainly. (i) \emph{Occlusion is an inspectable state}: named observed-versus-occluded regions a reader can enumerate, rather than a sub-symbolic visibility grid feeding a learned outcome model~\cite{curtis2024partially}; the limitation is that the visibility signal is currently supplied by an oracle rather than a perception module. (ii) \emph{Uncertainty is handled without probabilities}: Know-If Fluents record what is known deterministically, with no per-action probability to estimate; the limitation is that the planner cannot rank actions by success likelihood and instead minimises plan cost and replans when an optimistic assumption fails. (iii) \emph{No MDP machinery is required}: a lightweight deterministic PDDL layer replaces the learned-MDP stack of POMDP-based TAMP~\cite{curtis2024partially}; the limitation is that optimistic determinise-and-replan is sound only while belief reflects observation.

